Learning in artificial life requires experience to alter a continuing body's useful responses, yet acquired state, responsive physiology and material renewal can produce similar outward behaviour. We address this ambiguity with a comparative experimental framework that specifies the information available to a learner, intervenes on acquired content at matched bodily states and measures the resulting physical costs. Across distinct cellular models, local scalar consequences establish an action direction without a supplied reward gradient; controlled cue encounters support preference reversal and reassignment of temporal credit; and conditional experience changes the allocation of protective reserves. These effects have different limits. Scalar response estimation acquires in both tested conditions but reverses reward preference in only one, while a fixed responsive control also succeeds after an execution change. In a regenerative host, acquisition coincides with substantial tissue production, but newborn-content interventions do not establish that copying learned values is necessary for subsequent behaviour. At equal protection purchases, acquired conditional readout reduces injury by 16.6\% and 5.1\% relative to pooled content in two parents; its advantage over immediate physiology depends on environmental stability and contact timing. Together, the experiments distinguish useful acquired content from capabilities supplied by sensing, actuation and physiology. They establish conditional physical functions of experience, while leaving a unified, continuously adapting and regenerating learner unresolved.
